\chapter{Data Acquisition and Processing Pipeline}
\label{chap:data_pipeline}
% This is the heavy engineering chapter highlighting the 80% of the work that was data prep!

\section{Sourcing Transit Data}
\label{sec:sourcing_data}
The foundational step of this optimization problem was acquiring a comprehensive, time-accurate dataset of scheduled public transportation. The Guinness World Record guidelines mandate the use of scheduled public transport. Therefore, the data pipeline required a fusion of short-haul flight networks and high-speed rail schedules.

\medskip

\noindent The geographic scope of the dataset was strictly and intentionally limited to the European continent. Europe is the only logical theater for this specific record due to its unparalleled density of sovereign states coupled with highly developed, cross-border transit infrastructure. However, extracting continental-scale transit data presents a significant financial barrier. Commercial aggregators typically charge per query, which is prohibitively expensive when building a dense, multi-country matrix. Consequently, a primary constraint of this project was to secure massive data volumes while strictly minimizing or eliminating acquisition costs. This budget constraint fundamentally dictated the technological approach for both flight and train data and imposed some challenges for acquiring high quality data.

\subsection{Flight APIs}
\label{subsec:flight_apis}
While the airline industry is highly standardized, access to global distribution systems (such as Amadeus or Sabre) requires expensive commercial licensing. To bypass these costs, the flight dataset was sourced using the AviationEdge API\cite{aviationedge}.

\medskip

This platform offered a generous 1 month trial period charged at \$7 that enabled access to most of the endpoints, among those the \textit{flightsHistory} endpoint used to acquire the flight data. To efficiently target the API without exhausting query limits on irrelevant locations, an online global dataset of airports\cite{ourairports} was acquired and programmatically filtered to isolate only active commercial airports within European borders. The resulting list of IATA (International Air Transport Association) codes served as the exact querying parameters.

\medskip

By iterating through these European IATA codes and a specific 48-hour target date, the pipeline sent automated HTTP requests to extract all scheduled departures. The 48-hour window is required to ensure that routes starting at any point on day 1 have a full 24-hour window ahead of them to complete the record.

\subsection{Train Data}
\label{subsec:train_data}
Acquiring European rail data proved substantially more difficult than acquiring aviation data. Unlike the airline sector, which is strongly standardized around global distribution systems, European rail operations are fragmented across many national carriers (e.g., SNCF, DB, Renfe, and Trenitalia). No free, centralized API provides complete cross-operator schedules, while commercial rail datasets are prohibitively expensive for this project's budget.

\medskip

As a result, the rail pipeline relied on programmatic extraction from the \textit{Deutsche Bahn (DB) portal}\cite{deutschebahn}, which serves as a practical aggregator for a large share of international European rail connections. Because the DB interface is implemented as a modern Single-Page Application (SPA), data retrieval required browser-level automation rather than simple HTTP scraping.

\medskip

\noindent\textbf{Bot evasion and deterministic interaction.}
The scraper was implemented with \textit{undetected-chromedriver} in a headless Selenium workflow. Standard Selenium configurations are commonly flagged by web application firewalls through automation fingerprints (e.g., \texttt{navigator.webdriver}). To reduce detection risk, the browser session was initialized with Chrome DevTools Protocol overrides that masked automation indicators and emulated normal browser characteristics.

\medskip

The DB search UI also showed unstable behavior under keystroke-based automation because of dynamic autocomplete components and transient consent overlays. To improve robustness, form fields were populated through targeted JavaScript DOM injection (\texttt{driver.execute\_script}), followed by explicit event dispatches (\texttt{change} and \texttt{blur}) so that the SPA state logic updated consistently.

\medskip

\noindent\textbf{Incremental time-sweeping for full-day coverage.}
Na\"ive pagination via ``Later Connections'' frequently triggered asynchronous DOM replacements and stale element references. To avoid these failures, the scraper used an incremental time-sweeping strategy: after parsing one result page, it captured the departure time of the last listed service, relaunched the query with that time as the new lower bound, and repeated this process until the full 24-hour window was covered.

\medskip

\noindent\textbf{Bi-directional corridor extraction and checkpointing.}
To constrain the immense scale of the European rail network and prevent the scraper from blindly querying thousands of unviable, low-speed rural combinations, the extraction was strictly limited to a predefined compilation of major international rail corridors\cite{showmethejourney} sourced from train focused sites from the internet. This curated list of high-speed routes (e.g., London to Paris, Munich to Budapest) was hardcoded into the pipeline to serve as the exact boundaries of the dataset. The extraction algorithm iterated through this specific list, parsed the strings containing interconnected stops, isolated the terminal origin and destination nodes, and executed the incremental scraper in both forward and reverse directions.

\medskip

Given the computational cost and operational fragility of large-scale SPA extraction, the pipeline collected one representative weekday of rail schedules and duplicated that 24-hour block to form the 48-hour planning horizon required by the Guinness World Record routing formulation.

\section{Data Cleaning and Normalization}
\label{sec:cleaning_normalization}
With the raw flight and train schedules acquired, the data existed in highly disparate, semi-structured formats across thousands of CSV and JSON files. To safely process and mutate this data at scale, a local SQLite database was instantiated as an intermediate staging environment. This approach enabled structured SQL queries to filter, group, and reconcile the massive datasets before exporting the finalized, clean records back to CSV format for the graph construction algorithm.

\subsection{Entity Resolution and Standardization}
\label{subsec:entity_resolution_standardization}
A primary challenge in aggregating multi-modal European transit data is the absence of consistent naming conventions across transport networks. In practice, this meant that before a single route could be evaluated, the algorithm had to be taught that "München Hbf", "Munich Central", and "München Hauptbahnhof" were all the same physical place, and that "Frankfurt Flughafen" and "Frankfurt (Main) Hbf", despite sharing a city name, were emphatically not.

\medskip

\noindent\textbf{Toponymic normalization and multi-modal hub separation.}
To resolve this issue, a programmatic normalization pipeline mapped string variants to canonical entities. The algorithm first lowercased station text and then applied targeted heuristic matching to classify infrastructure types.

\medskip

Because the optimization problem is explicitly multi-modal, normalization also had to distinguish transport modes within the same metropolitan area. For example, variants of ``Frankfurt'' were explicitly separated into the canonical rail hub ``Frankfurt (Main) Hbf'' and the aviation node ``Frankfurt Flughafen''. Equivalent bifurcation rules were applied in Berlin (e.g., ``Berlin Hbf'' and ``Berlin S\"udkreuz'' versus ``Flughafen BER'') and Munich.

\medskip

The pipeline further unified strategically important cross-border junctions affected by strong linguistic variation. For example, heterogeneous labels around Basel and Lindau were merged into canonical station entities to preserve graph continuity at international interfaces.

\medskip

\noindent\textbf{Geopolitical entity standardization.}
Country labels also required strict harmonization between datasets. The aviation data provided ISO alpha-2 codes (e.g., ``GB'', ``MD''), which were mapped to textual country names using the \texttt{pycountry} library. However, these outputs did not always align with the naming used in rail records. Manual overrides were therefore introduced to prevent artificial country mismatches during route evaluation (e.g., mapping ``United Kingdom'' to ``England'', and simplifying ``Moldova, Republic of'' to ``Moldova'').

\medskip

\noindent\textbf{Deterministic temporal normalization.}
Although the flight API included historical and estimated delay attributes, these fields were removed from the final modeling dataset. Incorporating stochastic delay terms into a time-expanded graph substantially increases state-space complexity while offering limited predictive value for future schedules. The final normalized dataset therefore retained only scheduled departure and arrival times, preserving a deterministic Directed Acyclic Graph (DAG) formulation.

\subsection{The Midnight Rollover Problem and Absolute Minutes}
\label{subsec:midnight_rollover_absolute_minutes}
To allow mathematical evaluation of layover times, all chronological data had to be converted into a continuous timeline. The target metric was \textit{Absolute Minutes}, representing the elapsed time from $t=0$ (midnight on the first day of the 48-hour routing window).

\medskip

While flight records included explicit datetime strings, making baseline conversion straightforward, the scraped rail data provided mostly localized \textit{HH:MM} values for intermediate stops. Because rail itineraries frequently cross midnight, this created a \textit{Midnight Rollover} ambiguity: the raw strings did not indicate when the sequence advanced to the next calendar day, as the scraped data included only the times, not the dates. These were hardcoded in the scraping script, allowing us to inject them in the data manually and saving us from scraping another hard-to-extract field.

\medskip

To resolve this, a sequential tracking algorithm reconstructed absolute time along each route. For each event, local time was converted to minute-of-day form:
\[
M = \textit{hours} \times 60 + \textit{minutes}, \quad 0 \le M < 1440.
\]
Let $t_{prev}$ denote the previously validated absolute timestamp and $M_{current}$ the current local minute value. The previous minute-of-day is:
\[
M_{prev} = t_{prev} \bmod 1440.
\]
A rollover flag is raised when the local clock wraps around:
\[
R =
\begin{cases}
1, & \text{if } M_{current} < M_{prev}, \\
0, & \text{otherwise}.
\end{cases}
\]
With day index $D = \lfloor t_{prev} / 1440 \rfloor$, the current absolute timestamp is computed as:
\[
t_{abs} = (D + R)\times 1440 + M_{current}.
\]
This formulation unfolds cyclic 24-hour clock values into a strictly increasing timeline suitable for graph operations. In parallel, a conservative imputation rule was applied to partial records: when arrival or departure values were missing, the available counterpart was copied to avoid \texttt{NULL}-driven discontinuities in the routing graph.

\subsection{UTC Synchronization}
\label{subsec:utc_synchronization}
Because Europe spans multiple time zones, locally normalized $t_{abs}$ values are not directly comparable across borders. Without this step, the algorithm would perceive a train departing Madrid at 10:00 and a flight departing Lisbon at the same clock time as simultaneous events, when in reality the flight departs an hour earlier in absolute terms. A single unresolved timezone offset anywhere in the graph would silently corrupt every downstream routing decision that crossed that border.

\medskip

To ensure strict chronological consistency, every absolute timestamp was projected onto a UTC+0 reference axis. Let $\Delta_{offset}$ denote the local UTC offset in minutes. The synchronized value is:
\[
t_{UTC} = t_{abs} - \Delta_{offset}.
\]
Determining $\Delta_{offset}$ robustly across heterogeneous nodes required a cascading lookup strategy:

\begin{itemize}
  \item \textbf{Node-level overrides:} Explicit hardcoded mappings were applied first for exceptional nodes where country defaults are unreliable (e.g., Canary Islands $\Delta_{offset}=0$, Azores $\Delta_{offset}=-60$, and eastern Russian hubs such as Yekaterinburg $\Delta_{offset}=300$).
  \item \textbf{ISO-based country mapping:} If no node override existed, the ISO alpha-2 code was mapped to predefined regional offsets (e.g., Great Britain, Ireland, and Iceland to $0$; Finland, Greece, and Moldova to $120$).
  \item \textbf{Continental default:} Remaining nodes fell back to the Central European baseline $\Delta_{offset}=60$.
\end{itemize}
This hierarchy provided a stable and computationally efficient timezone normalization layer for the full dataset.

\medskip

After synchronization, a final integrity pass validated physical temporal feasibility for every leg:
\[
t_{UTC}^{dep} < t_{UTC}^{arr}.
\]
Any record violating this constraint was permanently removed, filtering upstream API inconsistencies and corrupted scrape artifacts from the final model.

\section{Spatial Mapping and Filtering}
\label{sec:spatial_mapping}
To evaluate a Guinness World Record attempt based on the number of countries visited, the routing algorithm must inherently ``know'' the geographical location and sovereign territory of every single node in the network. Furthermore, a raw European transit graph contains tens of thousands of minor stops (e.g., suburban commuter stations or rural villages) that drastically bloat the search space without offering viable international transfer opportunities. Consequently, a rigorous spatial mapping and filtering pipeline was engineered to reduce the graph to its most mathematically relevant hubs while ensuring 100\% geographical accuracy.

\subsection{Country Assignment and Border Resolution}
\label{subsec:country_assignment_border_resolution}
While international flights explicitly land in known countries, domestic and cross-border train schedules scraped from the DB portal do not consistently label the country of intermediate stops. To dynamically resolve this, a multi-strategy geographical detection algorithm was implemented\cite{kaggle_train_stations}.

\medskip

The algorithm evaluated each scraped station name through a descending hierarchy of spatial heuristics:

\begin{itemize}
  \item \textbf{Regex code extraction:} Extracting explicit country codes often hidden in parentheses within the raw text (e.g., parsing (CH) as Switzerland or (NL) as the Netherlands).
  \item \textbf{Language pattern recognition:} Utilizing linguistic markers to confidently infer sovereign territory (e.g., the presence of ``hl.n.'' for Czech main stations, ``Glowny'' for Poland, or `` b.Wien'' for Austrian suburbs).
  \item \textbf{GeoNames cross-referencing:} Querying a localized subset of the GeoNames database\cite{geonames} to map normalized station strings to the most populous European city matching that name.
\end{itemize}

Despite these programmatic layers, complex border stations frequently generated false positives (e.g., ``Gen\`eve'' being administratively in Switzerland despite a French linguistic profile, or border towns like Domodossola). To guarantee absolute integrity for record evaluation, an exhaustive list of manual topological corrections was hardcoded into the pipeline to forcefully overwrite erroneous geographical assignments at critical border crossings.

\subsection{Station Filtering and Hub Identification}
\label{subsec:station_filtering_hub_identification}
With sovereign territories assigned, the extracted graph still remained computationally intractable. The raw network contained many low-value commuter stops, rural halts, and suburban nodes that inflated branching factor without improving continental routing viability. A dedicated pruning pipeline was therefore applied to remove irrelevant nodes while preserving the international backbone.

\medskip

\noindent\textbf{Lexical and topological pruning.}
The first pruning layer used a deterministic station-scoring heuristic. Each node was evaluated by (i) topological frequency (how often it appeared across scraped routes) and (ii) lexical indicators of infrastructure scale. The scorer rewarded high-capacity terms (e.g., ``Hauptbahnhof'', ``Central'', ``Gare'', ``Aeropuerto'') and penalized low-capacity markers (e.g., ``Village'', ``Halt'', ``Suburban''). To prevent accidental graph fragmentation, terminal origin and destination stops of each train sequence were always retained, together with at least one high-scoring hub per traversed country.

\medskip

\noindent\textbf{LLM-assisted hub extraction.}
Although heuristic pruning removed most low-value nodes, identifying the most strategic international hubs across roughly 40 countries required additional geographic context. The pipeline therefore integrated Gemini 2.5 Pro as a constrained classifier rather than a routing engine.\cite{gemini2023} National station lists were provided after heuristic filtering, and the model was prompted to return only internationally valuable hubs (e.g., cross-border high-speed access and strong air--rail interconnectivity). Outputs were constrained to structured JSON for direct downstream use. This hybrid stage combined deterministic topological filtering with LLM-based geographic prioritization, yielding a compact graph concentrated around high-value transfer hubs.

\subsection{Intra-City Transfers and Multi-Agent Geocoding Validation}
\label{subsec:intra_country_transfers_geocoding_validation}
The final spatial challenge was the one that no algorithm could solve on its own: connecting otherwise separate transport modalities within metropolitan areas. Getting from a flight arriving at Paris CDG to a train departing Paris Gare du Nord is not a graph problem, it is a human logistics problem, with rush-hour traffic, terminal sprawl, and the very real possibility that the connection simply cannot be made in time. Every intra-city transfer link in the dataset needed a travel time that a real person, carrying luggage, could actually execute. Getting this wrong would not produce a suboptimal route. It would produce a route that fails in the middle of the attempt.

\medskip

\noindent\textbf{Multi-agent LLM architecture.}
Directly estimating transfer feasibility through commercial geocoding APIs at continental scale was cost-prohibitive. To bypass this constraint, a two-agent workflow was implemented:

\begin{itemize}
  \item \textbf{Generator agent:} evaluated hub pairs within the same country, proposed plausible intra-city connections, and estimated transit durations (local transit plus walking). A hard cutoff removed proposals with $\tau_{transfer} > 240$ minutes.
  \item \textbf{Auditor agent:} independently reviewed generated links, flagging implausible pairings and likely underestimates caused by geospatial hallucinations or naming ambiguities.
\end{itemize}

\medskip

\noindent\textbf{Ground-truth manual validation.}
The multi-agent stage reduced the candidate space to exactly 236 high-value intra-city transfer links. However, because LLM estimates remain probabilistic, a final manual validation pass was required for deterministic modeling quality. Each of the 236 links was manually verified in Google Maps, and exact walking/local-transit durations were extracted and hardcoded into the dataset. This final step eliminated subscription API costs while converting AI-proposed candidates into a physically viable, fully deterministic transfer matrix for routing optimization.

